\documentclass{article}

\usepackage{amsmath}
\usepackage{amssymb}
\usepackage{amsfonts}
\usepackage{listings}
\usepackage{xcolor}
\usepackage{graphicx}

\title{MS5612: Real Options Valuation}
\author{Lokesh Parihar BE22B008}
\date{\today}

\begin{document}
\maketitle
\newpage
\tableofcontents
\newpage

\section{27th July - Monday}
Finance could broadly be divided into 3 parts:
1. Persional Finance, Corporate Finance, and Financial Markets

How do businesses manages their financing? - Corp Fin is mostly about it.

Corp Fin could again broadly be classified into two: Investement and Financing.

Investement decisions result in assets and financing decisions result in liabilities.

Stocks could be valued in any of the following methods:
\begin{itemize}
    \item Comparative method
    \item Dividend Discount method
    \item Free Cash Flow method
    \item Book Value
\end{itemize}

How do we value Debt?
\begin{itemize}
    \item present value of cashflows of the debt (interest + principal)
\end{itemize}

How do we value projects (assets)?
\begin{itemize}
    \item Discounted Cashflow Method (DCM)
\end{itemize}

Keyword: Fortune 500
Keyword: Managers \textbf{create} shareholder value.

Most Fortune 500 firms use real options. But as per McKinzie's India head, almost no Indian company is using real options.
So, there is an opportunity.

\subsection{DCF vs Options Appoach}
\begin{itemize}
    \item DCF involves limited discretionary power, but options provide flexibilty.
    \item If people are liking flexibilty, we may value this flexibilty.
    \item Managers create flexibilty where seemingly none exists.
\end{itemize}


\section{29th July - Wednesday}

Focus of the course is on the management of long term assets (not much on liabilities).

NPV, IRR, Profitability Index (PI), Payback Period.

NPV - How much does the project worth? (\textbf{How much?})
\[NPV = \sum_{i = 1}^{n} \frac{C_i}{(1+r)^i} - IO\]

Keyword: IO - Initial Outlay

IRR - What returns are we making? (\textbf{What percent?})

PI - Present value of cash inflow/ initial outlay (IO) (\textbf{What rank?})

PB - \textbf{How long?}

\subsection{NPV} 

\[NPV = \sum_{i = 1}^{n} \frac{C_i}{(1+r)^i} - IO\]

Investment rule:  if $NPV > 0$ we may accept \\
if $NPV < 0$ we must reject \\
if $NPV = 0$ we reject it as it doesn't create any shareholder value.\\

\subsection{Class Discussion on Strategic Investment}
What are Strategic Investements?

How do you justify an Investement? NPV, IRR, PB, PI? But only these, really?

Some investments could have negative NPV, but they could still be chosen.

Strategic investments are those that may fail the test of conventional investment tests (NPV, IRR) but they may pass the tests in different cases.

\subsection{Example}
Let's say there's an RnD project requiring 15 million USD. 

To manufacture the plant, the could be 40 million, 80 million, and 120 million.

And the market share could be 50 million, or 130 million

The expected cost is 80 million, expected benefit is 90 million, and 15 millions for rnd.

Here, conventionally, the recommendation is to not investment and the NPV is -5 million USD.

\subsubsection{Case 1: Cost Certain, Revenue uncertain}
Let's say once we've invested in RnD, the manufacturing cost is certainly 80 million. If, now the market is 50 million, we do not invest. If however the market is 130 million, we go ahead and invest, in which case the benefit is 50 million. NPV in this case is \(\frac{1}{2} (0+15) - 15 = 10\) million. 


\subsubsection{Case 2: Revenue certain at 90, cost uncertain}
 if cost is 40 mil, we invest, benefit = 50 mil \\
 if cost is 80 mil, we invest, benefit = 10 mil \\
 if cost is 120 mil, we do not invest, benefit = 0 \\

 \(NPV = \frac{1}{3}(50 + 10 + 0) - 15 = 5 \) million.


 \subsubsection{Case 3: Revenue and cost are uncertain}
 (40, 50) = 10
 (40, 130) = 90
 (80, 50) = 0
 (80, 130) = 50
 (120, 50) = 0
 (120, 130) = 10

 NPV = 1/6(10 + 90 + 50 + 10) - 15 = 11.67 million

 In the formula of NPV, r is the measure of risk.

 The discount rate for a project must reflect the risk of the project.

 Case 3 is definitely the riskiest, but what about the other cases? Which one is more risky?
 We can calculate the risk? We may also calculate the variance. We'll notice that the NPV increases as the risk increases.

So, this RnD investment is an option that would give the firm the flexibilty to go ahead and set up a manufacturing plant. And the riskier the option is, the more valuable it is.


\section{03rd Aug - Monday}

\subsection{Summary of the last class}
\begin{itemize}
    \item Corporates have to only invest in projects having NPV positive, else they won't be creating shareholder values.
    \item But often there are Strategic investments.
    \item SIs may have negative NPV using the conventional DCF method.
    \item Strategic investments are those that have several options and if we add the value of these options, we may have positive NPV.
    \item As the uncertainty increases the value of the initial investments is higher.
\end{itemize}

\subsection{Options}
\begin{itemize}
    \item Options give the holder the right but not the obligation to do someting.
    \item \textbf{Call Option:} Provides the right to buy the underlying
    \item \textbf{Put Option:} Provides the right to sell the underlying
    \item These are called derivatives because they derive their value from the underlying.
    \item \textbf{Option Value:} Premium paid to buy the option/ market price of the option.
    \item \textbf{Option Holder:} Person/entity who's the buyer of the option.
    \item \textbf{Option Seller:} Person/entity who's selling the option.
\end{itemize}

\subsection{Options Payoff}
\begin{itemize}
    \item Stock - Simply a straight line with slope 1 against the value of Stock
    \item Risk free Bond - Simply a horizontal line against the value of Stock
    \item Call Option  
    \subitem For the buyer - $y = 0$ until $x = \text{strike}$, then a straight line slope 1. We may have to shift the graph downward by the amount paid as premium to get the profit.
    \subitem For the Seller - reflection of what it is for the buyer along the x-axis.
    \item The upside to the seller is limited but the downside is unlimited. Vice versa for the buyer.
    \item Options are a zero-sum game. Someone's gain is someone's loss.
    \item Put Option
    \subitem For the buyer - Reflection along $x = \text{strike}$ of the above curves
    \subitem For teh seller 
    \item Limited upsides and limited downsides
    \item Call and Put options are vanilla options. The fun comes when we combine some options.
\end{itemize}

\subsection{Popular Combinations of Options}
\begin{itemize}
    \item \textbf{Stock + Sold a Call}: Do from the graph
    \item \textbf{Risk Free Bond + Sold a Put}: Do from the graph
    \item The payoff of both the above are same.
    \item We do this exercise and can build a basket of options/securities to get the desired payoff.
\end{itemize}

\subsection{Let's get a little more creative}
\begin{itemize}
    \item \textbf{Buying a Call + Buying a Put:}  Almost a Kite, we make money, with high chance and the more volatile the stock is, the more chance.
    \item If we factor in the premium, there is a part on the x-axis for which the profit would be negative.
    \item \textbf{Long Call @ 100 + Short 2 Calls @ 110 + Long Call @ 120}: Butterfly Combination
    \item 
\end{itemize}

\subsection{Some real world examples of options}
\begin{itemize}
    \item Uber rides cost a little higher for reservation rides.
    \item Indigo charges higher for the option of low price cancellation. (Indigo Fare vs Indigo UpFront)
    \item What are the factors that influence this premium charge for cancellation?
    \item The premium seems to be the same across flights, accounting for flight distance, and time.
    \item Cancelling is a Put option.
    \item These are the examples of real options.
    \item Is insurance a put option?
\end{itemize}


\subsection{Course Project}
\begin{itemize}
    \item Step 1: Propose 3 three topics that we'd like to study
    \item Write a 100 words description on all three topics, indicating the use of real otpions.
    \item Submit hard copy by 5 pm, Aug 7th.
    \item Primary data and not desk based research.
\end{itemize}
























\end{document}
